\documentclass[a4paper,12pt]{article}
\usepackage[utf8]{inputenc}
\usepackage[T1]{fontenc}
\usepackage[margin=2.5cm]{geometry}
\usepackage{ctex} % 中文支持
\usepackage{xcolor}
\usepackage{titlesec}
\usepackage{setspace}

% -----------------------------------------------------------
% 格式设置区：Level 2 标准段落 + Level 3 法律/学术语态 + Level 2 合规披露
% -----------------------------------------------------------

% 设置层级标题格式（自动编号）
% 使用黑色，强调合规报告的严肃性
\titleformat{\section}{\Large\bfseries\color{black}}{\thesection}{1em}{}
\titleformat{\subsection}{\large\bfseries\color{black}}{\thesubsection}{1em}{}

% 设置行间距与首行缩进，适应长段落阅读
\onehalfspacing
\setlength{\parskip}{0.5em}
\setlength{\parindent}{2em}

% -----------------------------------------------------------
% 文档元数据
% -----------------------------------------------------------
\title{\textbf{2024年度环境、社会及管治（ESG）合规管理与绩效评价报告}}
\author{本公司董事会办公室}
\date{2025年3月}

\begin{document}

\maketitle

% 前言
\section*{前言}

本文件旨在依据适用的法律法规、行业监管规则及国际可持续发展报告准则，对本公司及其附属机构在2024财政年度内的环境管理体系运行状况、社会责任履行合规性及公司治理架构的有效性进行披露。文中陈述之数据均基于经审计的内部控制记录及第三方监测报告，旨在客观展示公司在可持续发展领域的合规绩效、风险管控能力及管理承诺的履行情况。

% 第一章 环境
\section{环境管理体系运行与合规性评价}

报告期内，本公司依据《环境保护法》及相关实施细则，维持并优化了环境管理体系（EMS）。该体系的运行旨在确保所有污染物排放均处于法定限值之内，并推动资源利用效率的持续提升，以满足日益严格的外部监管要求。

\subsection{环境合规绩效指标综述}

经核准的年度环境监测数据表明，公司在温室气体管控及废弃物合规处置方面达成既定管理目标。在排放管控方面，温室气体排放总量（涵盖范围1及范围2）被核定为39,107.65吨CO2e，经折算后的单位营收排放强度为4.58吨/百万元营收；一般工业固体废物产生量记录为170.48吨；对于具有环境风险的危险废物，其合规处置率维持在100\%水平，证明了所有有害废弃物均依据法定程序进行了无害化处理。在资源配置合规性方面，新鲜水取用总量为149.55万吨，包装材料消耗总量为7,677吨。为保障上述环境管理体系的有效运行及法律义务的履行，本年度环境保护专项资金投入总额为457.93万元，同时依法缴纳环境保护税9.6万元，体现了企业对环境外部成本的内部化承担。

\subsection{污染防治与清洁生产合规改造}

为响应国家关于工业领域碳达峰的政策导向，并确保生产设施符合最新的环保排放标准，一系列针对固定源及工艺流程的清洁化改造措施已被实施。在核心生产基地，依据《大气污染防治法》关于燃煤设施整治的要求，原有的燃煤供热系统已被清洁能源热源系统所替代。该工程的完工标志着该基地生产环节废气排放的完全消除，确保了大气污染物排放的绝对合规。与此同时，旗下化工厂已部署冷凝水回收利用装置。该装置通过物理回收生产蒸汽冷凝水并回注于锅炉补水系统，实现了水资源与热能的双重循环，此举被认定符合《节约用水条例》关于工业用水重复利用的规定。

针对西北某大型露天煤矿项目的移动源排放合规性问题，运输设备的电气化替代方案已被确立为核心管理策略。具体而言，21台纯电动矿用卡车及36台油电增程式矿用卡车被引入运营体系。上述设备的投入使用，旨在通过制动能量回收机制降低化石燃料消耗，从而从源头上减少移动源污染物的排放。此外，12台无人驾驶电动矿卡的应用不仅实现了运行阶段的零尾气排放，同时通过优化行驶路径降低了噪声污染，符合《环境噪声污染防治法》的相关要求。

% 第二章 创新
\section{技术创新体系与行业发展贡献}

技术创新被视作推动行业高质量发展及提升本质安全水平的关键驱动力。本公司持续投入研发资源，以构建具有自主知识产权的技术壁垒，并确保技术应用符合行业安全标准。

\subsection{研发合规投入与知识产权管理}

2024年度的研发活动资源配置及产出情况经审计确认符合高新技术企业认定的相关标准。研发支出总额为4.19亿元，该金额占当年营业收入的比例被核定为4.91\%。从事研发活动的在册专业人员总数为1,064人。在知识产权积累方面，本年度获得授权的专利数量为126项，截至报告期末，累计有效专利持有量为596项。上述知识产权的取得不仅确立了公司的技术所有权，也为公司的持续经营提供了法律层面的技术壁垒保护。

\subsection{核心技术推广与安全效益}

本年度重点推广的核心技术旨在解决行业共性技术难题并提升作业过程的安全性，符合国家关于“科技兴安”的政策导向。首先，“现场混装水胶炸药制备工艺”的实施，通过采用原材料分装运输、作业现场即时混合的流程，从根本上规避了成品炸药在长途运输中的公共安全风险，符合民用爆炸物品安全管理条例的导向；其技术填补了国内空白，密度可调特性允许根据地质条件进行能量匹配，提升了爆破作业的精细化水平。其次，“二氧化碳相变膨胀致裂技术”利用液态二氧化碳物理相变产生的高压气体破岩，全过程无化学燃烧与明火产生。该技术在国内首创并已形成500万只产能，被认定为适用于高瓦斯矿井等易燃易爆环境的优选安全技术方案。最后，“智能化矿山协同作业系统”基于5G通讯、边缘计算与自动化控制技术的深度融合，实现了钻机、检测机器人与运输车辆的无人化协同。该管理模式通过物理隔离将操作人员移出高危作业区域，显著降低了职业安全风险敞口。

% 第三章 社会
\section{社会责任履行与利益相关方权益保障}

本公司依据相关法律法规及社会公德准则，积极履行对员工、社区及社会的责任，确保企业运营产生的社会外部性为正向。

\subsection{职业健康安全（OHS）管理绩效}

截至2024年12月31日，公司劳动合同签订率维持在100\%水平，确保了劳动关系的合规性。为规避及转移运营风险，年度安全保障资金投入包括支付法定的安全生产责任保险费355.99万元，以及履行社会保险义务缴纳工伤保险费892.5万元。此外，公司建立了“1+14+N”应急预案体系以强化应急响应能力。在海外运营层面，纳米比亚矿业服务项目部实施了标准化的班前会议制度及“手指口述”安全确认程序，以确保跨文化背景下的作业规范性。报告期内，该项目达成连续130万工时无损失工时伤害（LTI）的记录，验证了国际化运营中安全管理体系的有效性及对国际劳工保护标准的遵循。

\subsection{企业公民义务与应急响应}

2024年度，公司对外捐赠及公益性支出总额为1,219万元，其中用于乡村振兴项目的专项资金为1,125万元，符合国家关于企业参与第三次分配的倡导。在突发公共事件应对方面，针对2024年湖南某地堤坝决口封堵行动，公司受命调动专业工程力量。作业团队采用了“松动爆破”控制技术，在严格控制震动参数的前提下，完成石料开采8.4万吨，为抢险工程提供了关键物资保障。该行动体现了企业在突发公共事件中的应急响应能力及对社会公共利益的维护。

% 第四章 治理
\section{公司治理架构与合规运营}

作为上级集团的成员企业，本公司遵循上市公司治理准则，维持了既定的治理架构与供应链管控标准，以保障投资者权益。

\subsection{治理结构有效性与信披质量}

董事会席位设置遵循“3+3+3”原则，即股东代表董事3名、独立董事3名及执行董事3名，该结构旨在确保决策权的制衡与独立性，防止利益冲突。在合规性评价方面，在证券交易所组织的年度信息披露质量评价中，本公司被授予“A级”评级，证明了公司信息披露的及时性、准确性与完整性。

\subsection{股东回报与商业道德管理}

依据公司章程及股东大会决议，2024年度实施了现金红利派发方案，总额为2.54亿元（每10股派发2.05元），该金额占当年可分配利润的40.12\%，切实保障了中小股东的收益权。在供应链合规管理方面，针对在库的1,548家供应商，公司强制执行了《廉洁协议》签署制度。经核查，协议签署率达到100\%，该措施旨在构建防御商业贿赂的合规防火墙，维护公平竞争的市场秩序，确保供应链运营的廉洁性。

\end{document}